%% =====================================================================
%%  B3-T-05-02 -- what B3 contributes to "Follow the Interest"
%%  -------------------------------------------------------------------
%%  LAYER A: APPEND-ONLY. Written once, then left alone. Material found
%%  only here sits in the `unique` slot and is protected by rule R10.
%% =====================================================================
%% @kind: source
%% @id: B3-T-05-02
%% @book: B3
%% @topic: T-05-02
%% @locator: Law 13, pp. 95--101; Law 7
%% @status: distilled
%% @unique: yes
%% @integrated: yes
%% @updated: 2026-09-19

\dossier{B3}{T-05-02}{\bookshort{B3}}{Law 13, pp. 95--101; Law 7}

\begin{distilled}
  \item Appeal to self-interest rather than to mercy or gratitude, because interest is something a person can act on immediately.
  \item Uncover what in your request benefits the other party and emphasise it out of proportion; they respond when they see something to gain.
  \item Using other people's work while taking the credit is presented as an economy of effort --- never do yourself what others can do for you --- which is the same tracing of benefit applied to your own position.
\end{distilled}

\begin{unique}
  \item The instruction to over-emphasise the counterpart's gain out of proportion, i.e. interest-tracing used offensively as a framing device rather than only as a diagnostic.
\end{unique}
